%!TEX root = ../csuthesis_main.tex
\begin{acknowledgements} 


本科四年的路，从开学报到走到现在写完这篇论文，转眼就到了要落笔说谢谢的时刻。

最先要感谢的是赵荣昌老师。课题方向从最初的几条候选里反复讨论，最后落到闭式解持续学习这条线上；这中间经历过几次推翻重来。每次我拿出半成型的想法时，您都愿意花很长时间，一点一点和我捋清楚问题在哪里、应该怎样继续往下走。论文落笔阶段，您在方法叙述、实验摆放、引用规范等方面的细致打磨，也让我真正看到一篇合格的学术论文应当是什么样子。

我也要感谢戚建宇师兄和李锐师兄。多模态那条线上有不少卡住的地方，比如模型训不动、显存不够，很多问题都是向两位师兄一点点问出来的。写第四章前的实验设计，也和两位师兄讨论过许多次。许多看似零散的建议，最后都实实在在地帮助我把工作往前推进了一步。

谢谢中南大学。在这个校园里读完人生的本科四年，专业上的积累、和老师同学相处的时光、一次次课程学习与科研尝试，都发生在这里。四年前刚入学时，我还并不清楚自己会走向怎样的方向；而现在，我即将从这里出发，保研直博至清华大学计算机系，继续在计算机科学与人工智能相关方向上学习和探索。这既是本科阶段努力的一个阶段性结果，也是一段新的旅程的开始。感谢中南大学给予我的培养、平台与机会，让我能够带着这四年的收获，走向下一段更广阔的求学道路。

最后想说的话，留给家人和朋友。论文里写不出来的那些事，比如投稿被拒后的低落，截稿前熬夜调实验时的疲惫，很多时候你们都一直在。你们的理解、支持和陪伴，是我能走到现在最大的底气。本科四年到这里即将画上句号，但这一路上得到的帮助、鼓励和温暖，我会一直记得。

\end{acknowledgements}
